% ============================ 一、问题重述 ============================
\section{问题重述}

\subsection{问题背景}

啤酒的风味是原料、酵母代谢与陈化过程共同作用的产物，其感官表现由数百种挥发性
与非挥发性化合物在极宽浓度区间内的协同与拮抗决定。酿造学界长期用“气味活度值
（odour activity value, OAV）”刻画单一化合物对香气的贡献，即
$\mathrm{OAV}=C/T$，其中 $C$ 为实测浓度、$T$ 为该化合物在啤酒基质中的气味阈值；
$\mathrm{OAV}>1$ 通常意味着该物质可被感知\cite{grosch1994,grosch1993}。然而
\textbf{真实啤酒是数百种化合物共存的多组分体系}：酯类贡献果香、萜烯类贡献酒花
草本香、高级醇贡献醇厚感、含硫化合物与醛酮类在微量时即构成缺陷。单一化合物的
OAV 无法解释“同样是 $\mathrm{OAV}<1$ 的两种物质，混合后香气显著增强”这类
协同效应，也难以回答“哪些化合物真正决定了某一条感官维度的评分”。

本题提供的商用啤酒风味代谢数据集，把这一难题转化为一个可计算的多视图建模问题：
既有 250 款商业啤酒在 231 个化学指标上的实测丰度，也有经专业品评小组按 0--6
分制评定的 50 个感官维度，还给出了化合物之间、感官维度之间、化学与感官之间的
三张相关矩阵，以及气味阈值与化合物来源标注。数据同时\textbf{真实地含有缺失值、
检出限以下截断与量纲混用}，这正是真实代谢组数据的常态。

\subsection{需要解决的问题}

题目要求围绕“化合物 $\rightarrow$ 感官”这一映射，完成三项递进的任务。

\textbf{问题一：数据缺失条件下的感官评分。}
以 \texttt{1.csv} 的化合物丰度为输入，给出 \texttt{3.xlsx} 所定义的 Aroma 与
Flavour 两套共约 40 个感官维度（连同口感、缺陷、香料与总体维度共 50 维）的
0--6 分制分级评分结果。难点在于：样本量仅 250 而特征 231 维，$\text{特征数}
\approx \text{样本数}$；感官标签本身是品评小组的均值，噪声大、部分维度极难预测；
化合物矩阵还存在系统性缺失与检出限截断。

\textbf{问题二：化合物--风味分析。}
利用 \texttt{2.csv} 的化合物相关矩阵构建化合物共表达图，结合 \texttt{1.csv}
的特征，为每个感官维度给出 Top-3 关键因素化合物，并对结果做\textbf{稳定性}
与\textbf{化学合理性}两类论证。稳定性需用 bootstrap 重采样下的入选频率、或不同
建模方法间 Top-3 的 Jaccard 系数来量化；化学合理性需要把统计关联还原到化合物
类别、代谢来源与气味阈值层面。

\textbf{问题三：风味与“风味指纹”生成。}
以 \texttt{tasting\_category\_fine} 为风格标签，学习\textbf{风格条件生成模型}：
在给定风格（如 Tripel、Kriek、Lambic）下生成合理的化合物谱，并给出对应的感官
画像；同时完成\textbf{跨风格迁移}，例如把 Pils/Lager 的化合物谱迁移为 Tripel
风格，且要求缺陷化合物（DMS、双乙酰、4VG、反式-2-壬烯醛、陈酒花等超阈值项）
\textbf{不恶化}。
